\partdivider{C}{Work, Energy and Collisions}{Two conserved quantities, energy and momentum, let you skip the details of a motion and jump straight from its start to its end. They outlive even the laws that produced them.}{\chref{ch:energy} Work, Energy and Power\\[1.2mm]
\chref{ch:collisions} Collisions}

\ncertchapter{11}{Work, Energy and Power}\label{ch:energy}

Newton's second law answers every question in mechanics in principle, but often only at great cost. To
find how fast a sled, starting from rest, is moving at the bottom of a curved hill by integrating $\vF=m\va$, one would need the
shape of the hill at every point and the direction of the normal force at every instant. Yet the answer, if
friction is negligible, depends only on the height of the hill. Something is being tracked that does not
care about the details of the path, and that something is energy.

This chapter introduces the first of the two great conservation principles of mechanics. It defines
work, the product of a force and a displacement along it, and proves from Newton's second law the
work--energy theorem: the net work done on a body equals the change in its kinetic energy, so the whole
history of a motion is traded for its two endpoints. When the forces are conservative, their work can be
stored as potential energy, and the sum of kinetic and potential energy stays constant.

The limits deserve stating plainly, because the method is so convenient that students use it where it
does not apply. When friction or drag does work, energy is not destroyed but leaves the mechanical
accounts as heat, and the equation must include it. And energy methods give speeds, not directions or
times; for those, the second law is still needed.

\section{The Dot Product and Work}

The first quantity to define is \textbf{work}, which measures how much a force contributes to changing a
body's motion over a stretch of its path. When you push a heavy box across a floor, only the part of your
force that lies \emph{along} the floor moves it forward; any part you direct downward or upward
contributes nothing to the motion along the floor. So work should count only the component of force along
the displacement, and that is exactly what the dot product of \chref{ch:vectors} extracts: for vectors
$\A$ and $\B$ of magnitudes $A$ and $B$ with angle $\theta$ between them,
$\A\cdot\B=AB\cos\theta=A_xB_x+A_yB_y+A_zB_z$, a scalar that reads ``how much of $\A$ lies along $\B$,
times the length of $\B$.'' The $\cos\theta$ is the projection that discards the perpendicular part.

\begin{defbox}{Work}
The work done by a constant force $\vF$ over a displacement $\vc{d}$ is
\[
W=\vF\cdot\vc{d}=Fd\cos\theta.
\]
For a varying force, split the path into small displacements $d\vr$ (with $\vr$ the position vector), over each of which $\vF$ is effectively
constant, and add the pieces $dW=\vF\cdot d\vr$: $W=\displaystyle\int \vF\cdot d\vr$. Work is a scalar,
measured in joules ($\SI{1}{J}=\SI{1}{N.m}$), and it carries a sign: positive when the force has a
component along the displacement, negative when it has a component against it. Physically, work is the
energy a force transfers: positive work feeds energy into the body, negative work takes it out
(Figure~\ref{fig:workangle}). For a force that varies along a straight line, the integral is an area
(Figure~\ref{fig:workarea}).
\end{defbox}

\begin{figure}[htbp]
\centering
\tikzfig{work_angle}
\figcaption{fig:workangle}{The angle between force and displacement sets the sign of the work. A force with a component along the motion ($\theta<90^\circ$) does positive work; one with a component against it ($\theta>90^\circ$) does negative work; a perpendicular force ($\theta=90^\circ$) does \emph{no} work. If it is the only force doing work, it therefore speeds the body up, slows it down, or leaves its speed unchanged.}
\end{figure}

\begin{trapbox}{a force can act and do no work}
The normal force on a block sliding over a fixed surface, the tension holding a mass in circular motion,
gravity on a horizontally moving object: all act at $90^\circ$ to the motion and do \emph{zero} work. ``A force is
present'' and ``work is done'' are different claims, though such a force is not idle: tension in circular
motion turns the velocity without changing the speed. Nor is work effort: holding a heavy bag still tires

your arm, but the bag does not move, so your force on it does no work.
\end{trapbox}

\begin{figure}[htbp]
\centering
\tikzfig{work_area}
\figcaption{fig:workarea}{For a variable force, work is the area under the force--displacement graph: each thin strip contributes $dW=F\,dx$, and the strips add up to $W=\int F\,dx$. For a spring gun (spring constant $k$, initial compression $x_0$) the force falls linearly from its maximum $kx_0$ to zero as the spring relaxes, so the area is a triangle, $W=\tfrac12kx_0^2$.}
\end{figure}

\begin{exbox}{Work from the dot product}
With $\uhat{\imath}$ and $\uhat{\jmath}$ the unit vectors along $x$ and $y$,
a force $\vF=(6\,\uhat{\imath}+8\,\uhat{\jmath})\ \si{N}$ acts on a crate as it undergoes a displacement
$\vc{d}=(4\,\uhat{\imath}-3\,\uhat{\jmath})\ \si{m}$. Find the work done, and the angle between force and
displacement.

\stepn{1} \textbf{Work is the dot product.} Both vectors are given in components, so use the component form:
\[
W=\vF\cdot\vc{d}=(6)(4)+(8)(-3)=24-24=\SI{0}{J}.
\]
\stepn{2} \textbf{Interpret.} Zero work with a nonzero force and displacement means they are
\emph{perpendicular}. Check: $|\vF|=\sqrt{6^2+8^2}=10$, $|\vc{d}|=\sqrt{4^2+3^2}=5$, and
$\cos\theta=W/(|\vF||\vc{d}|)=0$, so $\theta=90^\circ$. On its own, this force would not change the
crate's speed.
\ans{$W=0$, $\theta=90^\circ$}
\end{exbox}

\section{Kinetic Energy and the Work--Energy Theorem}

Work has been defined; what does it actually \emph{do} to a body? It changes one particular quantity,
which should be named first. The \textbf{kinetic energy} of a body of mass $m$
moving at speed $v=|\vv|$ is
\[
K=\tfrac12 m v^2=\tfrac12 m\,(\vv\cdot\vv),
\]
a scalar (joules), never negative (zero only at rest), and, being built from $\vv\cdot\vv$, independent of the
direction of motion. In terms of the momentum $p=mv$ of \chref{ch:newton}, $K=p^2/2m$. The theorem below
shows that $K$ is the work needed to bring the body from rest to speed $v$: the energy of motion. Like
speed, $K$ depends on the frame of reference.

\begin{derivbox}{The work--energy theorem}
We connect the net force to the change in kinetic energy, with no reference to time, keeping everything
a vector until the dot product turns it into a scalar.

\stepn{1} \textbf{Start from Newton's second law, as a vector.}
\[
\vF_{\text{net}}=m\,\frac{d\vv}{dt}.
\]

\stepn{2} \textbf{Form the work done over an element of the path.} The work by the net force over a small
displacement $d\vr$ is the dot product $dW=\vF_{\text{net}}\cdot d\vr$. Since velocity is $\vv=d\vr/dt$,
the step taken in time $dt$ is $d\vr=\vv\,dt$, so
\[
dW=\vF_{\text{net}}\cdot d\vr=m\,\frac{d\vv}{dt}\cdot\vv\,dt=m\,(\vv\cdot d\vv).
\]

\stepn{3} \textbf{Recognise the dot product as a perfect differential.} Since
$\dfrac{d}{dt}(\vv\cdot\vv)=2\,\vv\cdot\dfrac{d\vv}{dt}$, we have $\vv\cdot d\vv=\tfrac12\,d(\vv\cdot\vv)
=\tfrac12\,d(v^2)$. Therefore
\[
dW=\tfrac12 m\,d(v^2)=d\!\left(\tfrac12 m v^2\right)=dK.
\]
Each little bit of work is exactly the little change in kinetic energy it produces.

\stepn{4} \textbf{Integrate over the whole path,} from the start point (speed $u$) to the end (speed $v$);
the prime on $v'$ marks the integration variable, which runs from $u$ to $v$:
\[
W_{\text{net}}=\int \vF_{\text{net}}\cdot d\vr=\int_{u}^{v} d\!\left(\tfrac12 m v'^2\right)
=\tfrac12 m v^2-\tfrac12 m u^2=\Delta K.
\]
The dot product projected force onto motion at every step, so we never had to assume the motion was
along the force.
\end{derivbox}

\begin{keybox}{Work--energy theorem}
\[
\boxed{\;W_{\text{net}}=\int\vF_{\text{net}}\cdot d\vr=\Delta K=\tfrac12 mv^2-\tfrac12 mu^2\;}
\]
It holds for \emph{any} force, constant or varying, conservative or not, in one dimension or three,
because it is nothing but Newton's second law dotted with the displacement and integrated. It applies to
a single particle, or to a body whose parts all move together; for a system whose parts move differently,
internal forces can also do work (a compressed spring pushing two blocks apart gives both of them kinetic
energy, though no external force acts). It needs an inertial frame; in an accelerating frame the
pseudo-force of \chref{ch:circdyn} must be included as one more force doing work. (In 1-D along $x$
the dot product reduces to $W=\int F\,dx$; for a \emph{constant} force this is
$Fs=\tfrac12mv^2-\tfrac12mu^2$, which is the $v^2=u^2+2as$ of \chref{ch:straightline} in energy
language.)
\end{keybox}

\inwords{The theorem replaces ``track the force vector through every instant'' with ``compare the two
speeds,'' which is why energy methods are so much faster whenever you neither know nor want the time.}

\section{Power}

The work--energy theorem says nothing about time, yet for an engine time is what matters: the stronger
engine delivers the same energy sooner. Work says how much energy is transferred; power says how quickly.

\begin{defbox}{Power}
\textbf{Power} is the rate of doing work. If a force $\vF$ acts on a body moving with velocity $\vv$, the
instantaneous power it delivers is
\[
P=\frac{dW}{dt}=\vF\cdot\vv,
\]
since $dW=\vF\cdot d\vr$, and dividing by $dt$ turns $d\vr$ into $\vv$. Again it is a dot product: only
the component of force along the velocity delivers power. The average power
over a time $t$ is $\bar P=W/t$ (the bar means ``average''). Power is a scalar, measured in watts
($\SI{1}{W}=\SI{1}{J/s}$; do not confuse the upright unit W with the italic symbol $W$ for work). When
force and velocity are parallel (a car engine driving straight ahead), the dot product reduces to $P=Fv$.
\end{defbox}

$P=Fv$ also explains why a car has a top speed. An engine of fixed power $P$ pushes with $F=P/v$, which
shrinks as the car speeds up (\SI{60}{kW} gives \SI{6000}{N} at \SI{10}{m/s}, \SI{2000}{N} at
\SI{30}{m/s}), while drag grows (Figure~\ref{fig:power}). Acceleration stops when $F$ equals the total
resistance $F_{\text{res}}$: $v_{\max}=P/F_{\text{res}}$. The example below uses $P=\vF\cdot\vv$ the
other way round, to find the power a known force demands.

\begin{exbox}{The power to climb}
A cyclist of mass \SI{60}{kg} (bicycle included) rides up a slope at a steady \SI{4}{m/s}, gaining
height at \SI{1.2}{m/s} (the vertical component of her velocity). What power does she supply against
gravity ($g=\SI{9.8}{m/s^2}$)?

\stepn{1} \textbf{Power of gravity.} We do not know her pedalling force, but we know gravity, and her
power must balance it. With $y$ upward, gravity is $\vF=-mg\,\uhat{\jmath}$ with
$mg=60\times9.8=\SI{588}{N}$, so in $P=\vF\cdot\vv$ only the vertical component $v_{\text{vert}}$ of her
velocity survives: $P_{\text{grav}}=-mg\,v_{\text{vert}}=-588\times1.2\approx\SI{-706}{W}$. Neither the
slope's angle nor her \SI{4}{m/s} speed along it is needed.
\stepn{2} \textbf{Her power.} Her speed is steady, so her kinetic energy does not change and the net
power is zero. Ignoring friction and drag, she must therefore supply $+\SI{706}{W}$, about one
horsepower (\SI{746}{W}), which a trained rider can sustain only briefly. \ans{$\approx\SI{0.7}{kW}$}
\end{exbox}

\begin{figure}[htbp]
\centering
\tikzfig{power_cubic}
\figcaption{fig:power}{Why top speed is expensive. Aerodynamic drag rises as $v^2$, and power is force times speed, so the power to overcome drag rises as $v^3$. Doubling a car's top speed takes roughly \emph{eight} times the engine power.}
\end{figure}

\section{Conservative Forces and Potential Energy}

The work--energy theorem still needs the work of every force along the actual path. For some forces, though, the work done between two points does not depend on the route taken.
For such a force the work could be worked out once, from the endpoints alone.
Gravity near the Earth is the standard example. Take $y$
upward, so $\vF=-mg\,\uhat{\jmath}$; along any
small step $d\vr=dx\,\uhat{\imath}+dy\,\uhat{\jmath}$ the work is $dW=\vF\cdot d\vr=-mg\,dy$, because the
horizontal part of the step contributes nothing. Adding the steps of \emph{any} path from height $y_1$
to $y_2$ gives $W=-mg(y_2-y_1)$: a straight drop and a curved slide between the same heights get the same work. Forces like this allow work to be stored and
recovered, and the stored quantity is potential energy.

\begin{defbox}{Conservative force}
A force is \emph{conservative} if the work it does depends only on the endpoints, not the path (equivalently,
the work around any closed loop is zero, since out by one route and back by another gives zero exactly
when the two routes give equal work).
Gravity and the spring force are conservative; friction is not.
For a conservative force we can define a \emph{potential energy} $U$ (a scalar, in joules) by
$W_c=-\Delta U$, where $W_c$ is the work done by the conservative force: that work equals the drop in
$U$. Equivalently, $\Delta U$ is the work an
external agent must do to move the body slowly (with no change of kinetic energy) against the force. Over
a small step $dx$ this reads $F\,dx=-dU$, so in one dimension $F=-\dfrac{dU}{dx}$ (in three dimensions,
each component is obtained the same way: $F_x=-\partial U/\partial x$, where $\partial/\partial x$ means
differentiate with respect to $x$ holding $y$ and $z$ fixed, and likewise for $y$ and $z$).
\end{defbox}

\begin{remarkbox}{Why friction is not conservative}
The work done by kinetic friction depends on the \emph{path length}, not the endpoints. Slide a book in a big loop back to where it started: gravity has done zero net work (same
height), but friction has opposed the motion \emph{every centimetre} of the trip, draining energy the
whole way. A force whose work depends on how far you
travelled, not just where you ended up, cannot be summarised by a potential energy $U(x)$, because
$U$ is a function of \emph{position} alone. This is the same distance-versus-displacement split from
\chref{ch:straightline}, now with energy consequences.
\end{remarkbox}

\begin{trapbox}{potential energy belongs to a pair, not to one body}
``The ball's potential energy'' is shorthand. $U=mgh$ is the energy of the Earth--ball \emph{system}, set by
their separation; a spring's $\tfrac12kx^2$ is stored in the spring, not in the block attached to it. The
shorthand is harmless when one partner (the Earth, a wall) barely moves, so that all the
kinetic-energy change happens in the other body.
\end{trapbox}

\begin{figure}[htbp]
\centering
\tikzfig{pe_curve}
\figcaption{fig:pecurve}{The force is minus the slope of the potential energy, so it always points ``downhill'' on this curve (arrows along the bottom): minima are stable equilibria and maxima unstable ones, with no calculation. Close to a minimum where the curvature is not zero, $U$ looks like $\tfrac12kx^2$ (dashed parabola), so a body near stable equilibrium behaves like a mass on a spring for small displacements. That is why oscillations turn up throughout physics.}
\end{figure}

\begin{keybox}{Standard potential energies}
Gravity near Earth: $U=mgh$, with $h$ the height above a chosen level. A spring of spring constant $k$
(in \si{N/m}, from Hooke's law in \chref{ch:connected}) stretched or compressed by $x$: $U=\tfrac12 kx^2$.
Universal gravitation: $U=-\dfrac{GMm}{r}$ for masses $M$ and $m$ a distance $r$ apart, with $G$ the
gravitational constant (\chref{ch:gravitation}). In each case $F=-dU/dx$ (with $x$ replaced by $h$ or $r$)
recovers the force. Only \emph{differences} of $U$ enter any physical result, so the zero is a choice:
$h=0$ at any convenient level, $x=0$ at the spring's natural length, and $U=0$ at infinite separation
for universal gravitation. The form $mgh$ assumes $g$ is constant, so it holds only for heights small
compared with the Earth's radius.
\end{keybox}

The gravity result is the path-independent work found at the start of this section, $W_c=-mg\Delta h$,
so $\Delta U=mg\Delta h$. The spring result takes one line. As the spring is stretched from its natural
length to $x$, the force it exerts at an intermediate extension $x'$ (primed, as before, to mark the
integration variable) is $-kx'$, so the work done by the spring is
$-\int_0^x kx'\,dx'=-\tfrac12kx^2$ (the triangle under the $F$--$x$ line). Since $W_c=-\Delta U$, the
stored energy is $U=\tfrac12kx^2$.

A graph of $U(x)$ shows the force at a glance (Figure~\ref{fig:pecurve}), because $F=-dU/dx$ is minus the
slope. Where the curve is flat, $dU/dx=0$, the force vanishes and the body can rest in
\textbf{equilibrium}. At a minimum ($d^2U/dx^2>0$) a small displacement either way meets a slope that
pushes the body back, so the equilibrium is \textbf{stable}; at a maximum ($d^2U/dx^2<0$) the slope
pushes it further away, so the equilibrium is \textbf{unstable}. The curve also shows where motion is
allowed. Call $E=K+U$ the body's \textbf{mechanical energy}; if it stays constant (the subject of the
next section), then since $K=E-U$ cannot be negative, the body is confined to the region
where $U(x)\le E$, and it turns round where $U=E$.

\section{Conservation of Mechanical Energy}

Potential energy was invented so that the work of conservative forces could be booked in advance. Putting
it into the work--energy theorem now gives the most used result of the chapter: when only conservative
forces do work, kinetic energy and potential energy trade with each other and their sum, the mechanical
energy $E=K+U$, stays fixed. Figure~\ref{fig:energybars} shows that trade during a free fall.

\begin{figure}[htbp]
\centering
\tikzfig{energy_bars}
\figcaption{fig:energybars}{Energy bookkeeping during a free fall from height $H$. The potential-energy bar (purple)
shrinks exactly as fast as the kinetic-energy bar (orange) grows, so the total column height, the
mechanical energy $mgH$, is identical at every stage.}

\end{figure}

\begin{derivbox}{When mechanical energy is conserved}
Every force acting on a body is either conservative (gravity, springs) or not (friction, drag, tension,
applied pushes). Split the net work into the work $W_c$ of the conservative forces and the work
$W_{\text{nc}}$ of the non-conservative ones, $W_{\text{net}}=W_c+W_{\text{nc}}$. The work--energy theorem
says the net work is the change in kinetic energy:
\[
W_c+W_{\text{nc}}=\Delta K.
\]
\stepn{1} \textbf{Use the definition of potential energy.} Potential energy is \emph{defined} so that the
work done by a conservative force equals the drop in potential energy: $W_c=-\Delta U$. (For gravity,
lifting a body does negative gravitational work and \emph{raises} $U=mgh$; the signs match.)
\stepn{2} \textbf{Substitute and regroup.} Replacing $W_c$ by $-\Delta U$:
\[
-\Delta U+W_{\text{nc}}=\Delta K\ \Rightarrow\ W_{\text{nc}}=\Delta K+\Delta U=\Delta(K+U).
\]
\stepn{3} \textbf{Read off the conclusion.} The non-conservative work equals the change in total
mechanical energy. So if \emph{no} non-conservative force does work
 ($W_{\text{nc}}=0$), then $K+U$ does
not change; it is conserved.
\end{derivbox}

\begin{keybox}{Conservation of mechanical energy}
When only conservative forces do work, the total mechanical energy is constant:
\[
K+U=\tfrac12 mv^2+U=\text{constant}
\]
(Figure~\ref{fig:energycons} shows the trade for a pendulum).

In general $\Delta(K+U)=W_{\text{nc}}$. For friction from a fixed surface or drag from still air,
$W_{\text{nc}}<0$ and the lost energy appears as heat; an applied push or a motor-driven tension can make
$W_{\text{nc}}>0$. (Friction can do \emph{positive} work on one body, as when static friction
carries a crate forward on an accelerating lorry; but kinetic friction always removes mechanical energy
from the two surfaces taken together.) The condition is that non-conservative forces do no \emph{work},
not that they are absent: the normal force on a smooth roller-coaster track acts throughout but is always
perpendicular to the motion, so mechanical energy is still conserved.
\end{keybox}

\begin{trapbox}{three different ``conservations''}
Total energy, heat included, is \emph{always} conserved; mechanical energy $K+U$ only when
$W_{\text{nc}}=0$; kinetic energy alone hardly ever (in a free fall it grows as $U$ falls).
\end{trapbox}

\begin{figure}[htbp]
\centering
\tikzfig{energy_pendulum}
\figcaption{fig:energycons}{A pendulum trades potential for kinetic energy and back: all potential at the extremes (momentarily at rest), all kinetic at the bottom (maximum speed), and half of each at half the height, with the sum fixed throughout. The same bookkeeping handles roller-coasters, ski-jumps and dropped balls.}
\end{figure}

\begin{exbox}{Drop and friction}
A block of mass $m$ is released from rest at height $h$ on a smooth curved track. At the bottom it runs
onto a rough horizontal floor with coefficient of kinetic friction $\mu$ (the $\mu_k$ of
\chref{ch:friction}). How far, $d$, does it slide before stopping?

The two stages obey different rules, so treat them separately: on the track only gravity does work; on
the floor friction does.

\stepn{1} \textbf{Smooth descent (energy conserved).} From top to the bottom of the drop,
$mgh=\tfrac12 mv^2$, so the block arrives with $v=\sqrt{2gh}$ and kinetic energy $K=mgh$.
\stepn{2} \textbf{Rough floor (energy removed by friction).} On the level floor the normal force is
$N=mg$, so the friction force is $\mu mg$, opposing the motion; over distance $d$ it does work
$W_{\text{nc}}=-\mu mg\,d$. The block stops, so $\Delta K=0-mgh=-mgh$, and $U$ does not change on the
level floor. Setting $W_{\text{nc}}=\Delta K$,
\[
-\mu mg\,d=-mgh\quad\Longrightarrow\quad d=\frac{h}{\mu}.
\]
\stepn{3} \textbf{Sanity check.} Mass cancels: a heavy block and a light one stop after the same
distance, because both the driving energy ($mgh$) and the friction force ($\mu mg$) scale with $m$.
\ans{$d=h/\mu$}
\end{exbox}

\section{Summary}\label{sec:summary11}
Work counts only the component of force along the motion; the work--energy theorem turns Newton's second
law into one scalar equation between endpoints; and a force whose work depends only on the endpoints has
a potential energy, which makes $K+U$ a conserved quantity when only such forces do work.

\begin{itemize}
\item $W=\vF\cdot\vc{d}=Fd\cos\theta$; for a varying force, $W=\int\vF\cdot d\vr$. A perpendicular force
      does no work.
\item A conservative force has $W_c=-\Delta U$ and $F=-dU/dx$.
\item $W_{\text{net}}=\Delta K=\tfrac12mv^2-\tfrac12mu^2$, valid for every kind of force.
\item Power $P=dW/dt=\vF\cdot\vv$; at top speed the driving force equals the resistance.
\item Only differences of $U$ matter; the zero of potential energy is a choice.
\item Standard potential energies: $mgh$ near the Earth, $\tfrac12kx^2$ for a spring, and $-GMm/r$ outside a spherical
      mass, with $U=0$ at infinity (\chref{ch:gravitation}).
\item In general $\Delta(K+U)=W_{\text{nc}}$; for friction from a fixed surface or drag from still air,
      $W_{\text{nc}}<0$ and the missing mechanical energy appears as heat.
\item On a potential-energy curve, equilibrium sits where $dU/dx=0$; a minimum is stable and a maximum is
      unstable. A body of energy $E$ can only be where $U\le E$.
\end{itemize}

\section{Problem Types}\label{sec:types11}
Six types cover the work and energy questions examiners set, and most share one lesson: when a question
asks for a speed and the path is complicated, reach for energy before forces.

\archq{Energy-method problems}
{Speeds or heights wanted, with a force that varies or a time that is unknown.}
{Write $K+U$ at two points: $K_1+U_1+W_{\text{nc}}=K_2+U_2$, with $W_{\text{nc}}=0$ if no friction or drag acts. No
need to find accelerations or times.}
{Using energy conservation when friction is present without including its work.}
{A \SI{2}{kg} block slides from rest down a curved track of height \SI{5}{m} and reaches the bottom at
\SI{8}{m/s} ($g=\SI{10}{m/s^2}$). The work done on it by friction is:}
{\SI{-36}{J}}
{\SI{+36}{J}}
{\SI{-64}{J}}
{\SI{-100}{J}\archans{A}{$W_{\text{nc}}=\Delta K+\Delta U=\tfrac12(2)(8^2)-(2)(10)(5)=64-100=\SI{-36}{J}$.
Option (B) has the wrong sign: friction removes mechanical energy here; (C) and (D) mistake one of the
two energy terms for the whole change.}}

\archq{Work done by a variable or vector force}
{``Given $\vF$ (in components, or as a function of position) and a displacement, find the work,'' or
``find the work done by a spring / by gravity over a path.''}
{For a constant force use $W=\vF\cdot\vc{d}$ (components: $F_xd_x+F_yd_y+F_zd_z$). For a position-varying
force use $W=\int\vF\cdot d\vr$, the area under the $F$--$x$ curve in 1-D. Spring: the work done \emph{by the spring} as it is
stretched from its natural length by $x$ is $W=-\tfrac12 kx^2$; work done by gravity: $W=-mg\,\Delta h$.}
{Work is a scalar with a sign; a perpendicular force does zero work no matter how large. Do not multiply
magnitudes without the $\cos\theta$.}
{A force $\vF=(3\uhat{\imath}+4\uhat{\jmath})\,$N moves a body through
$\vc{d}=(2\uhat{\imath}+1\uhat{\jmath})\,$m. The work done is:}
{\SI{6}{J}}
{\SI{10}{J}}
{\SI{14}{J}}
{\SI{2}{J}\archans{B}{$W=\vF\cdot\vc{d}=(3)(2)+(4)(1)=6+4=\SI{10}{J}$. Option (A) keeps only the
$x$-terms; multiplying magnitudes, $|\vF||\vc{d}|=5\sqrt5\approx11.2$, ignores the angle.}}

\archq{Work--energy theorem for the net force}
{``A body's speed changes from $u$ to $v$ over a distance, find the net work, or the stopping
distance under a given force.''}
{$W_{\text{net}}=\Delta K=\tfrac12 m(v^2-u^2)$. For stopping under a constant force $F$,
$F\,d=\tfrac12 mu^2$ gives $d=\dfrac{mu^2}{2F}$. Works for \emph{any} force without finding time.}
{Net work uses the \emph{net} force (include friction). Stopping distance $\propto u^2$: double the speed,
quadruple the distance.}
{A \SI{2}{kg} body moving at \SI{6}{m/s} is brought to rest by a constant force in \SI{3}{m}. The
magnitude of the force is:}
{\SI{6}{N}}
{\SI{12}{N}}
{\SI{18}{N}}
{\SI{24}{N}\archans{B}{$W_{\text{net}}=\Delta K$: $-Fd=0-\tfrac12 mu^2$, so
$F=\dfrac{mu^2}{2d}=\dfrac{2\times36}{2\times3}=\SI{12}{N}$. Option (D) drops the $\tfrac12$ in the
kinetic energy.}}

\archq{Power problems}
{``Find the power of a pump/engine/person raising a mass, or the power at a given speed against a given
resistance.''}
{Use $P=\dfrac{W}{t}$ for an average, or $P=\vF\cdot\vv$ for an instant. For a pump lifting water,
include \emph{both} energy sinks if the water also leaves with speed: $P=\dot m\big(gh+\tfrac12v^2\big)$,
where $\dot m$ (the dot means ``per second'') is the mass of water lifted each second, in \si{kg/s}.
At top speed the driving force equals the resistance $F_{\text{res}}$, so $v_{\max}=P/F_{\text{res}}$.}
{Forgetting that at \emph{constant} speed the net work is zero: the power is spent entirely against
resistance or gravity, not on accelerating.}
{A \SI{60}{kg} person climbs a \SI{5}{m} staircase in \SI{10}{s} ($g=\SI{10}{m/s^2}$). The power
developed against gravity is:}
{\SI{30}{W}}{\SI{120}{W}}{\SI{300}{W}}{\SI{3000}{W}\archans{C}{$P=\dfrac{mgh}{t}
=\dfrac{60(10)(5)}{10}=\SI{300}{W}$. Option (D) is the work, $mgh=\SI{3000}{J}$, not yet divided by
the time.}}

\archq{Spring energy and the energy method}
{``A block compresses a spring / is launched by one, find the compression, or the speed at some
point.''}
{Equate energies between two instants: $\tfrac12kx^2$ for the spring, $\tfrac12mv^2$ for motion,
$mgh$ for height, and subtract $\mu mg\,d$ if friction acts. Never use $F=kx$ with a
constant-acceleration formula; the force varies.}
{Forgetting that a block descending onto a spring keeps falling \emph{while} the spring compresses, so
the height dropped is $(d+x)$, not $d$.}
{A spring of force constant \SI{500}{N/m} is compressed by \SI{20}{cm}. The energy stored is:}
{\SI{5}{J}}{\SI{10}{J}}{\SI{20}{J}}{\SI{50}{J}\archans{B}{$U=\tfrac12kx^2
=\tfrac12(500)(0.2)^2=\tfrac12(500)(0.04)=\SI{10}{J}$. Option (C) drops the $\tfrac12$ and (D) forgets
to square $x$; leaving $x$ in centimetres gives an absurd \SI{e5}{J}.}}

\archq{Potential-energy curve, equilibrium and stability}
{``Given $U(x)$, find the equilibrium positions and say whether each is stable,'' or ``find the force at
a point.''}
{Force is the negative slope: $F=-\dfrac{dU}{dx}$. Equilibrium is where $F=0$, that is,
$\dfrac{dU}{dx}=0$. Then use the second derivative: $\dfrac{d^2U}{dx^2}>0$ (a \emph{minimum}) is
\textbf{stable}; $<0$ (a \emph{maximum}) is \textbf{unstable}.}
{Reading stability off the sign of $U$ rather than its curvature. A minimum can sit at negative $U$ and
still be perfectly stable; what matters is the shape, not the value.}
{For $U(x)=2x^3-9x^2+12x$ (SI units), the equilibrium positions are:}
{$x=1$ only}{$x=2$ only}{$x=1$ (unstable) and $x=2$ (stable)}{$x=1$ (stable) and $x=2$
(unstable)\archans{C}{$\dfrac{dU}{dx}=6x^2-18x+12=6(x-1)(x-2)=0$ at $x=1,2$. Then
$\dfrac{d^2U}{dx^2}=12x-18$, which is $-6$ at $x=1$ (maximum $\Rightarrow$ unstable) and $+6$ at $x=2$
(minimum $\Rightarrow$ stable). Option (D) has the stabilities swapped, the error of reading a maximum
as a resting place.}}

